\documentclass[12pt]{article}

\usepackage{sbc-template}
\usepackage{graphicx,url}
\usepackage[utf8]{inputenc}
\usepackage[brazil]{babel}
\usepackage{amsmath}
\usepackage{booktabs}
\usepackage{multirow}

\sloppy

\title{Heurística Construtiva Gulosa para o Problema de Atribuição\\
Quadrática Generalizada (GQAP)}

\author{Yasmin Souza\inst{1}, José Guedes\inst{1}}

\address{Sistemas de Informação -- Universidade Federal de Viçosa (UFV)\\
  Campus Rio Paranaíba -- Rio Paranaíba, MG -- Brasil
}

\begin{document}

\maketitle

\begin{abstract}
The Generalized Quadratic Assignment Problem (GQAP) is an NP-Hard combinatorial
optimization problem that extends the classical Quadratic Assignment Problem by
incorporating capacity constraints. This paper presents a greedy constructive
heuristic for a simplified version of the GQAP, in which each entity is
iteratively assigned to the feasible location that minimizes the incremental
cost, considering both the linear allocation cost and the quadratic interaction
cost with already-assigned entities. Computational experiments were conducted on
45 synthetic instances of three sizes and three capacity tightness levels. All
instances of medium and large sizes were solved feasibly, while tight small
instances revealed limitations of the greedy approach -- motivating future
extensions toward metaheuristics such as GRASP.
\end{abstract}

\begin{resumo}
O Problema de Atribuição Quadrática Generalizada (GQAP) é um problema de
otimização combinatória NP-Difícil que estende o clássico QAP incorporando
restrições de capacidade. Este trabalho apresenta uma heurística construtiva
gulosa para uma versão simplificada do GQAP, na qual cada entidade é
iterativamente atribuída à localização viável que minimiza o custo incremental,
considerando o custo linear de alocação e o custo quadrático de interação com as
entidades já alocadas. Experimentos computacionais foram realizados em 45
instâncias sintéticas de três tamanhos e três níveis de apertamento de
capacidade. Todas as instâncias de médio e grande porte foram resolvidas de
forma viável, enquanto instâncias pequenas com capacidade apertada revelaram
limitações da abordagem gulosa, motivando extensões futuras com
meta-heurísticas como o GRASP.
\end{resumo}


% =============================================================================
\section{Introdução}
% =============================================================================

A Pesquisa Operacional (PO) é uma área interdisciplinar que aplica métodos
matemáticos e computacionais ao apoio à tomada de decisão em sistemas complexos.
Dentre os problemas tratados pela PO, os de \textit{atribuição} se destacam pela
sua ampla aplicabilidade: alocação de tarefas a máquinas, posicionamento de
instalações, design de redes de comunicação, entre outros.

O Problema de Atribuição Quadrática (QAP, do inglês \textit{Quadratic
Assignment Problem}) é um dos mais estudados nesse contexto e pertence à classe
dos problemas NP-Difíceis. Sua versão generalizada, o GQAP (\textit{Generalized
Quadratic Assignment Problem}), introduz restrições de capacidade nas
localizações, tornando o problema ainda mais desafiador para métodos exatos
\cite{cordeau:2006}.

\cite{mateus2011grasp}

Para problemas dessa natureza, heurísticas e meta-heurísticas são as abordagens
predominantes na literatura. As heurísticas construtivas, em particular,
desempenham papel fundamental como ponto de partida para algoritmos mais
sofisticados, como o GRASP (\textit{Greedy Randomized Adaptive Search
Procedure}) \cite{resende:2011}, que combina construção gulosa-aleatória com
busca local.

Este trabalho apresenta a \textbf{Etapa 1} de um projeto de pesquisa incremental:
a implementação de uma heurística construtiva gulosa para o GQAP simplificado.
O algoritmo serve como base para extensões futuras que incorporarão componentes
de aleatoriedade (RCL) e busca local, evoluindo para um GRASP completo nas
próximas etapas.

As contribuições deste trabalho são: (i) modelagem formal do GQAP simplificado,
(ii) implementação e análise de uma heurística construtiva gulosa com custo
incremental, (iii) geração de um conjunto de instâncias sintéticas reproduzíveis,
e (iv) experimentos computacionais com análise dos resultados.


% =============================================================================
\section{Definição do Problema}
% =============================================================================

\subsection{O Problema de Atribuição Quadrática Generalizada}

O GQAP pode ser descrito da seguinte forma: dado um conjunto de $n$ entidades
$E = \{1, \ldots, n\}$ e um conjunto de $m$ localizações $L = \{1, \ldots, m\}$,
deseja-se atribuir cada entidade a exatamente uma localização, respeitando as
capacidades das localizações, de forma a minimizar o custo total da atribuição.

\subsection{Formulação Matemática}

Seja $x_{ij} \in \{0,1\}$ a variável de decisão tal que $x_{ij} = 1$ se a
entidade $i$ é atribuída à localização $j$, e $x_{ij} = 0$ caso contrário.
A função objetivo do GQAP simplificado é:

\begin{equation}
\min \; F(x) = \sum_{i=1}^{n} \sum_{j=1}^{m} c_{ij} \, x_{ij}
             + \sum_{i=1}^{n} \sum_{k=i+1}^{n} \sum_{j=1}^{m} \sum_{l=1}^{m}
               f_{ik} \, d_{jl} \, x_{ij} \, x_{kl}
\label{eq:fo}
\end{equation}

\noindent sujeito a:

\begin{equation}
\sum_{j=1}^{m} x_{ij} = 1, \quad \forall \, i \in E
\label{eq:atrib}
\end{equation}

\begin{equation}
\sum_{i=1}^{n} w_{ij} \, x_{ij} \leq \text{cap}_j, \quad \forall \, j \in L
\label{eq:cap}
\end{equation}

\begin{equation}
x_{ij} \in \{0, 1\}, \quad \forall \, i \in E, \; j \in L
\label{eq:bin}
\end{equation}

onde:
\begin{itemize}
  \item $c_{ij}$: custo linear de alocar a entidade $i$ na localização $j$;
  \item $f_{ik}$: fluxo (interação) entre as entidades $i$ e $k$;
  \item $d_{jl}$: distância entre as localizações $j$ e $l$;
  \item $w_{ij}$: consumo de recursos da entidade $i$ na localização $j$;
  \item $\text{cap}_j$: capacidade total da localização $j$.
\end{itemize}

O primeiro somatório da Equação~\ref{eq:fo} representa o \textit{custo linear}
de alocação. O segundo representa o \textit{custo quadrático}, que penaliza a
co-alocação de entidades com alto fluxo mútuo em localizações distantes entre
si -- característica central do GQAP.

A restrição~\ref{eq:atrib} garante que cada entidade é atribuída a exatamente
uma localização. A restrição~\ref{eq:cap} impõe que o consumo total de recursos
em cada localização não exceda sua capacidade.


% =============================================================================
\section{Algoritmo Heurístico}
% =============================================================================

\subsection{Estratégia Construtiva Gulosa}

A heurística proposta constrói a solução incrementalmente: a cada passo, uma
entidade ainda não alocada é atribuída à localização que minimiza o
\textit{custo incremental}, definido como a variação no valor da função objetivo
ao inserir essa entidade.

Para a entidade $i$ sendo avaliada para a localização $j$, o custo incremental
$\Delta(i, j)$ é calculado como:

\begin{equation}
\Delta(i, j) = c_{ij} + \sum_{\substack{k \in E \\ x_k \neq -1}} f_{ik} \cdot d_{j, s(k)}
\label{eq:delta}
\end{equation}

onde $s(k)$ é a localização já atribuída à entidade $k$. A localização é
viável apenas se $w_{ij} \leq \text{cap}_j^{\text{restante}}$.

\subsection{Pseudocódigo}

\begin{verbatim}
Algoritmo: Heurística Construtiva Gulosa para o GQAP
Entrada: custo_linear, fluxo, distancia, consumo, capacidade
Saída: solucao[], custo_total

1.  Inicializar solucao[i] = -1  para todo i
2.  cap_restante[j] = capacidade[j]  para todo j
3.  Para cada entidade i de 0 a n-1:
4.      melhor_local  = -1
5.      melhor_custo  = +infinito
6.      Para cada localização j de 0 a m-1:
7.          Se consumo[i][j] > cap_restante[j]: continua
8.          delta = custo_linear[i][j]
9.          Para cada entidade k já alocada (solucao[k] != -1):
10.             delta += fluxo[i][k] * distancia[j][solucao[k]]
11.         Se delta < melhor_custo:
12.             melhor_custo = delta
13.             melhor_local = j
14.     Se melhor_local != -1:
15.         solucao[i] = melhor_local
16.         cap_restante[melhor_local] -= consumo[i][melhor_local]
17.     Senão:
18.         Registrar i como não alocada (solução inviável)
19. Retornar solucao, calcular_custo_total(solucao)
\end{verbatim}

\subsection{Estruturas de Dados}

O algoritmo utiliza as seguintes estruturas:

\begin{itemize}
  \item \textbf{Matrizes 2D} ($n \times m$) para \texttt{custo\_linear} e
        \texttt{consumo};
  \item \textbf{Matrizes 2D simétricas} ($n \times n$) para \texttt{fluxo} e
        ($m \times m$) para \texttt{distancia};
  \item \textbf{Vetor} \texttt{solucao[n]}: mapeia cada entidade à sua
        localização atribuída;
  \item \textbf{Vetor} \texttt{cap\_restante[m]}: rastreia a capacidade
        disponível em cada localização.
\end{itemize}

\subsection{Análise de Complexidade}

O algoritmo executa três laços aninhados:

\begin{itemize}
  \item Laço externo: $n$ entidades;
  \item Laço intermediário: $m$ localizações candidatas;
  \item Laço interno (custo quadrático): até $n$ entidades já alocadas.
\end{itemize}

A complexidade de tempo do algoritmo é portanto $O(n^2 \cdot m)$.
A complexidade de espaço é $O(n^2 + m^2 + nm)$, dominada pelas matrizes
de fluxo e distância.

Para as maiores instâncias testadas (30 entidades, 10 localizações), isso
corresponde a $30^2 \times 10 = 9.000$ operações elementares -- o que explica
os tempos de execução na ordem de décimos de milissegundo observados nos
experimentos.


% =============================================================================
\section{Experimentos Computacionais}
% =============================================================================

\subsection{Ambiente Computacional}

Os experimentos foram executados em um computador com processador Intel Core
i-série, sistema operacional Windows 11 e linguagem Python 3.x (sem uso de
bibliotecas de otimização externas). O código foi desenvolvido utilizando apenas
a biblioteca padrão do Python, garantindo reprodutibilidade e portabilidade.

\subsection{Instâncias Sintéticas}

Na ausência de datasets padronizados para a versão simplificada do GQAP adotada,
foram geradas 45 instâncias sintéticas por meio de um gerador controlado
com semente aleatória (\textit{seed}) para garantir reprodutibilidade.

As instâncias foram organizadas em três grupos de tamanho e três níveis de
\textit{apertamento de capacidade} (fator $\alpha$), totalizando 9 combinações
com 5 instâncias cada:

\begin{itemize}
  \item \textbf{Tamanhos}: pequeno ($5 \times 3$), médio ($15 \times 6$),
        grande ($30 \times 10$) -- notação entidades $\times$ localizações;
  \item \textbf{Fatores de capacidade}: folgado ($\alpha = 2{,}0$),
        médio ($\alpha = 1{,}5$), apertado ($\alpha = 1{,}1$).
\end{itemize}

Os valores de custo linear foram amostrados uniformemente em $[1, 100]$;
os fluxos em $[0, 50]$; as distâncias em $[1, 20]$; e os consumos em $[1, 10]$.
A capacidade de cada localização $j$ foi definida como:

\begin{equation}
\text{cap}_j = \alpha \cdot \bar{w}_j \cdot \frac{n}{m}
\end{equation}

onde $\bar{w}_j$ é o consumo médio das entidades na localização $j$.

\subsection{Resultados}

A Tabela~\ref{tab:resultados} apresenta os resultados agregados por grupo de
instâncias. Para cada combinação de tamanho e fator de capacidade são reportados:
número de instâncias viáveis, custo médio, custo mínimo, custo máximo e tempo
médio de execução em milissegundos.

\begin{table}[ht]
\centering
\caption{Resultados da heurística construtiva gulosa nas instâncias sintéticas
         (5 instâncias por grupo)}
\label{tab:resultados}
\begin{tabular}{llcrrrr}
\toprule
\textbf{Tamanho} & \textbf{Capacidade} & \textbf{Viáveis} &
\textbf{Custo Médio} & \textbf{Custo Mín.} & \textbf{Custo Máx.} &
\textbf{Tempo (ms)} \\
\midrule
\multirow{3}{*}{$30 \times 10$}
  & Folgada ($\alpha=2{,}0$)   & 5/5 & 61.464,03 & 48.932,57 & 82.675,86 & 0,6058 \\
  & Média ($\alpha=1{,}5$)     & 5/5 & 77.183,02 & 68.109,71 & 94.998,69 & 0,7575 \\
  & Apertada ($\alpha=1{,}1$)  & 5/5 & 85.833,44 & 74.697,77 & 95.585,15 & 0,6830 \\
\midrule
\multirow{3}{*}{$15 \times 6$}
  & Folgada ($\alpha=2{,}0$)   & 5/5 & 11.209,45 & 7.144,25  & 16.088,11 & 0,1105 \\
  & Média ($\alpha=1{,}5$)     & 5/5 & 17.042,48 & 12.374,32 & 24.588,39 & 0,1216 \\
  & Apertada ($\alpha=1{,}1$)  & 5/5 & 22.493,62 & 18.320,10 & 25.302,69 & 0,0792 \\
\midrule
\multirow{3}{*}{$5 \times 3$}
  & Folgada ($\alpha=2{,}0$)   & 5/5 & 1.133,68  & 685,91    & 1.884,12  & 0,0333 \\
  & Média ($\alpha=1{,}5$)     & 5/5 & 1.555,62  & 580,15    & 2.816,43  & 0,0284 \\
  & Apertada ($\alpha=1{,}1$)  & \textbf{2/5} & 1.459,81 & 779,45 & 2.934,29 & 0,1661 \\
\bottomrule
\end{tabular}
\end{table}

\subsection{Análise dos Resultados}

\textbf{Viabilidade.} A heurística encontrou soluções viáveis em 42 das 45
instâncias (93,3\%). As 3 instâncias inviáveis pertencem exclusivamente ao grupo
pequeno com capacidade apertada ($5 \times 3$, $\alpha = 1{,}1$). Nesse cenário,
a estratégia puramente gulosa pode esgotar a capacidade de certas localizações
nas primeiras alocações, inviabilizando as entidades restantes. Esse resultado
evidencia uma limitação conhecida das heurísticas construtivas gulosas:
\textit{decisões localmente ótimas podem comprometer a viabilidade global}.

\textbf{Qualidade das soluções.} O custo médio cresce significativamente com o
apertamento da capacidade. Para instâncias grandes, o custo médio na configuração
apertada (85.833) é aproximadamente 40\% superior ao da configuração folgada
(61.464). Essa diferença é esperada: com menos opções viáveis, a heurística
não consegue sempre escolher a localização de menor custo incremental.

\textbf{Desempenho computacional.} Todos os tempos de execução foram inferiores
a 1 ms, mesmo para as maiores instâncias ($30 \times 10$). Esse comportamento é
consistente com a complexidade teórica $O(n^2 \cdot m)$: para $n=30$ e $m=10$,
o número de operações é 9.000, trivial para hardware moderno. A escalabilidade
da heurística é, portanto, uma de suas principais vantagens.

\textbf{Variância entre sementes.} Em alguns grupos, observa-se variação
considerável entre o custo mínimo e máximo (por exemplo, 48.932 a 82.675 nas
instâncias grandes folgadas). Isso indica que a qualidade da solução é sensível
à estrutura particular de cada instância, o que reforça a necessidade de
mecanismos de diversificação -- como a aleatorização via RCL do GRASP -- nas
etapas futuras.


% =============================================================================
\section{Conclusão}
% =============================================================================

Este trabalho apresentou uma heurística construtiva gulosa para o GQAP
simplificado, com função objetivo composta por termos linear e quadrático, e
restrições de capacidade por localização. O algoritmo possui complexidade
$O(n^2 \cdot m)$ e demonstrou eficiência computacional excelente, resolvendo
instâncias de até 30 entidades e 10 localizações em menos de 1 ms.

Os experimentos revelaram que a heurística é confiável para instâncias com
capacidade folgada ou média, mas apresenta limitações em cenários de capacidade
muito restrita, onde a estratégia gulosa pode levar a soluções inviáveis.

Como trabalhos futuros, planeja-se:
\begin{itemize}
  \item \textbf{Etapa 2:} introdução da Lista Restrita de Candidatos (RCL) para
        aleatorizar a construção, transformando o algoritmo em um GRASP;
  \item \textbf{Etapa 3:} incorporação de uma fase de busca local
        (ex.: troca de pares, 2-opt) para refinamento das soluções construídas;
  \item \textbf{Etapa 4:} validação com instâncias da OR-Library e comparação
        direta com os resultados de \cite{cordeau:2006}.
\end{itemize}

\bibliographystyle{sbc}
\bibliography{referencia}

\end{document}
